% =====================================================================
% Report content to be pasted INSIDE \begin{document} ... \end{document}
%
% Required packages in your Overleaf preamble:
%   \usepackage{graphicx}
%   \usepackage{amsmath, amssymb}
%   \usepackage{float}        % for [H] figure placement
%   \usepackage{booktabs}     % nice tables (\toprule, \midrule, \bottomrule)
%   \usepackage{listings}     % code snippets
%   \usepackage{xcolor}
%   \usepackage{hyperref}     % cross-references
%   \lstset{basicstyle=\ttfamily\footnotesize, breaklines=true,
%           keywordstyle=\color{blue}, commentstyle=\color{gray!80!black},
%           numbers=left, numberstyle=\tiny, stepnumber=1, frame=single,
%           language=Matlab}
%
% Place the folder 'images_project3' next to your main.tex
% =====================================================================

\title{Adaptive Cruise Control Design \\ \large Project 3 -- Driver Assistance System Design}
\author{Dennis [Surname] \\ Politecnico di Torino}
\date{\today}
\maketitle

%======================================================================
\section{Introduction}
%======================================================================
Adaptive Cruise Control (ACC) is one of the most widespread longitudinal
driver-assistance functions: it autonomously regulates the speed of the
ego vehicle in order to keep a safe, time-dependent distance from a
detected preceding vehicle. The European standard
ISO~15622:2018 specifies the minimum set of Key Performance Indicators
(KPIs) that a Type~I ACC system must satisfy in terms of operating
speed, automatic acceleration and deceleration, and jerk. This project
designs and validates an ACC controller compliant with that standard
on a full nonlinear longitudinal vehicle simulator.

The vehicle plant is a complete powertrain model that accepts the
normalised accelerator pedal $h_a\in[0,1]$ and the normalised brake
pedal $h_b\in[0,1]$ as inputs and returns the time histories of the
longitudinal speed $v$, longitudinal acceleration $\dot v$ and travelled
distance $x$. Since the high-level ACC controller naturally generates
a target longitudinal acceleration $a_\text{target}$, a low-level layer
is needed to convert $a_\text{target}$ into the pair $(h_a,h_b)$. The
project is therefore organised in three steps:
\begin{itemize}
    \item \textbf{Step 1 -- Plant inversion only.} The simplified
          longitudinal dynamic equation is inverted analytically and
          implemented as a feed-forward function
          (\texttt{inverse\_pedals}). The first simulation uses this
          function as the sole low-level controller.
    \item \textbf{Step 2 -- Augmenting the inversion with a PI loop.}
          The plant-model mismatch causes a steady-state drift in the
          acceleration tracking. A PI controller is added on the
          acceleration error to recover zero-error tracking.
    \item \textbf{Step 3 -- Compliance verification.} The augmented
          ACC is tested against two distinct lead-vehicle speed
          profiles, one mild and one aggressive, in dry road
          conditions; the ISO~15622:2018 KPIs are enforced by a
          dedicated block in series with the ACC output.
\end{itemize}

Throughout the project the vehicle parameters of Genta Vol.\,2,
Appendix~E.1 (saloon configuration) are used. The most relevant ones
are summarised in Table~\ref{tab:params_p3}.

\begin{table}[H]
    \centering
    \begin{tabular}{lcl}
        \toprule
        \textbf{Symbol} & \textbf{Value} & \textbf{Description} \\
        \midrule
        $m$         & $1020$\,kg               & Vehicle mass \\
        $R_e$       & $0.279$\,m               & Rolling radius \\
        $T_{e,\max}$& $105$\,N\,m              & Max engine torque \\
        $T_{b,\max}$& $300$\,N\,m              & Max front braking torque \\
        $K_b$       & $2.3$                    & Brake bias ratio (70/30) \\
        $\eta_t$    & $0.93$                   & Driveline efficiency \\
        $\tau_f$    & $1/3.438$                & Final drive ratio \\
        $J_e,J_t,J_w$&$0.088,\,0.05,\,0.59$\,kg\,m$^2$ & Engine, trans., wheel inertia \\
        $S$         & $2.04$\,m$^2$            & Frontal area \\
        $C_x,C_z$   & $0.33,\,0.20$            & Drag / lift coefficients \\
        $f_0$       & $0.011$                  & Rolling resistance constant \\
        $K$         & $2.6\!\times\!10^{-8}$\,s$^2$/m$^2$ & Speed-dependent rolling term \\
        \bottomrule
    \end{tabular}
    \caption{Vehicle parameters used in Project~3 (saloon set).}
    \label{tab:params_p3}
\end{table}

%======================================================================
\section{Theoretical Background}
\label{sec:theory_p3}
%======================================================================

\subsection{Longitudinal vehicle model and its inversion}
The simplified longitudinal dynamics of the vehicle, lumped on a single
equivalent mass $m_e$, reads
\begin{equation}
    m_e\,\dot v =
        \underbrace{\frac{T_{e,\max}\,\eta_g\,\eta_t}{R_e\,\tau_g\,\tau_f}}_{F_{a,\max}(\tau_g)}\,h_a
        \;-\;
        \underbrace{\frac{T_{b,\max}}{R_e}\!\left(1+\tfrac{1}{K_b}\right)}_{F_{b,\max}}\,h_b
        \;-\;
        \bigl(A + B\,v^{2} + C\,v^{4}\bigr),
        \label{eq:long_dyn}
\end{equation}
where the equivalent mass
\begin{equation}
    m_e = m + \frac{4 J_w}{R_e^2} + \frac{J_t}{R_e^2 \tau_f^2} +
          \frac{J_e}{R_e^2\,\tau_g^2\,\tau_f^2}
    \label{eq:meq}
\end{equation}
accounts for the rotational inertia of wheels, transmission and
engine, and the rolling/aerodynamic resistance coefficients are
\begin{equation}
    A = m g\,f_0,\qquad
    B = m g K + \tfrac{1}{2}\rho S\,(C_x - C_z f_0),\qquad
    C = -\tfrac{1}{2}\rho S K C_z.
    \label{eq:resistance_p3}
\end{equation}

Equation~\eqref{eq:long_dyn} can be rewritten in compact form as
$\dot v = f(v,\tau_g,h_a,h_b)$. Imposing $\dot v\equiv a_\text{target}$
gives the inversion law that the low-level controller must implement:
the required wheel force is
\begin{equation}
    F_\text{req}(v,a_\text{target},\tau_g)
        = m_e\,a_\text{target} + A + B v^2 + C v^4,
    \label{eq:Freq}
\end{equation}
and a positive value is delivered through the throttle ($h_b=0$),
\begin{equation}
    h_a = \frac{F_\text{req}}{F_{a,\max}(\tau_g)}
        = \frac{R_e\,\tau_g\,\tau_f}{T_{e,\max}\,\eta_g\,\eta_t}\;
          F_\text{req},
    \label{eq:ha}
\end{equation}
whereas a negative value triggers the brake ($h_a=0$),
\begin{equation}
    h_b = \frac{-F_\text{req}}{F_{b,\max}}
        = \frac{-R_e\,F_\text{req}}{T_{b,\max}\bigl(1+1/K_b\bigr)}.
    \label{eq:hb}
\end{equation}
Because $h_a,h_b\in[0,1]$ are mutually exclusive (no throttle and brake
overlap is allowed) and the engine cannot push the wheels backwards,
the sign of $F_\text{req}$ uniquely selects the active pedal.

\subsection{Constant Time-Gap (CTG) policy}
The high-level ACC computes $a_\text{target}$ from the relative
distance and speed with respect to the preceding vehicle. The desired
inter-vehicle distance is proportional to the ego speed:
\begin{equation}
    L_\text{des}(v) = h\,v,
    \label{eq:Ldes}
\end{equation}
with $h$ the design time-gap (here $h=2.7$\,s). Defining the spacing
error
\begin{equation}
    \varepsilon = (x_l - x_e) - L_\text{des}(v_e)
                = (x_l - x_e) - h\,v_e,
    \label{eq:eps}
\end{equation}
its derivative $\dot\varepsilon = (v_l - v_e) - h\,a_e$ contains both
the velocity mismatch and the ego acceleration. The control law adopted
in the project is the classical CTG PD law
\begin{equation}
    a_\text{target} = -\frac{1}{h}\bigl(\lambda\,\varepsilon + \dot\varepsilon\bigr),
    \label{eq:CTG}
\end{equation}
where $\lambda = 0.8$ is the spacing-error gain. Equation~\eqref{eq:CTG}
combines proportional action on the position error with derivative
action on the velocity mismatch, and the prefactor $1/h$ guarantees the
correct dimensional scaling.

A minimum-set-speed switch is added before computing $L_\text{des}$:
if the ego is travelling below $v_\text{set,min}=7$\,m/s, the formula
uses $v_\text{set,min}$ instead of $v_e$, in compliance with the ISO
requirement on the minimum set speed.

\subsection{ISO 15622:2018 KPIs}
The standard prescribes five quantitative limits on the ACC behaviour,
listed in Table~\ref{tab:iso}. They are enforced downstream of the
CTG block (see Section~\ref{sec:simulink}).

\begin{table}[H]
    \centering
    \begin{tabular}{lc}
        \toprule
        \textbf{KPI} & \textbf{Limit} \\
        \midrule
        Min. operating speed for automatic acceleration & $v_\text{min,acc} = 5$\,m/s \\
        Minimum set speed                               & $v_\text{set,min} = 7$\,m/s \\
        Max. average deceleration (2\,s window)         & $3.5$\,m/s$^2$ \\
        Max. deceleration rate (negative jerk, 1\,s)    & $2.5$\,m/s$^3$ \\
        Max. automatic acceleration                     & $2.0$\,m/s$^2$ \\
        \bottomrule
    \end{tabular}
    \caption{Key Performance Indicators imposed by ISO~15622:2018.}
    \label{tab:iso}
\end{table}

%======================================================================
\section{Simulink Implementation}
\label{sec:simulink}
%======================================================================

\subsection{Overall architecture}
The full simulator is implemented in \texttt{DASD\_B\_Project\_3.slx}
and is illustrated in Figure~\ref{fig:sim_full}. The signal flow, from
left to right, is the following:
\begin{enumerate}
    \item the \textbf{Lead Vehicle} block plays back a speed profile
          and exports the leader bus $[x_l,\,v_l]$;
    \item the \textbf{ACC Controller} reads the leader bus together
          with the ego bus $[x_e,\,v_e]$ and computes
          $a_\text{target}$ according to the CTG law~\eqref{eq:CTG};
    \item the \textbf{ISO 15622:2018} block clamps and slope-limits
          $a_\text{target}$ to satisfy the standard KPIs;
    \item the PI controller (Section~\ref{sec:pi}) compensates for the
          residual acceleration error and feeds the inverse-pedals
          function;
    \item the \textbf{Ego Vehicle} block (the full powertrain
          simulator) integrates the longitudinal dynamics under the
          pedal commands $(h_a,h_b)$ and returns the ego bus.
\end{enumerate}
A stop condition triggers when the time gap becomes negative
(\texttt{<=0} block fed to a \texttt{STOP}), to flag any collision
between ego and leader.

\begin{figure}[H]
    \centering
    \includegraphics[width=\textwidth]{images_project3/simulink_full.png}
    \caption{Top-level architecture of the Project~3 Simulink model.}
    \label{fig:sim_full}
\end{figure}

\subsection{ACC controller block}
The internal structure of the ACC controller is shown in
Figure~\ref{fig:sim_acc}. The preceding bus is demultiplexed into
$x_l$ (sent to \texttt{out.x\_leader}) and $v_l$ (sent to
\texttt{out.v\_leader}); the ego bus is demultiplexed into $x_e$ and
$v_e$. The block in the top-right corner performs the minimum-set-speed
switch
$
    v_\star = \max\bigl(v_e,\,v_\text{set,min}\bigr),
$
multiplied by $h$ to produce the desired headway $L_\text{des}$. The
spacing error $\varepsilon = (x_l - x_e) - L_\text{des}$ is built by
two summing junctions, differentiated by a $\Delta u/\Delta t$ block
to obtain $\dot\varepsilon$, scaled by $\lambda$ and finally combined
through the $-1/h$ gain as prescribed by~\eqref{eq:CTG}. The time gap
$t_\text{gap} = (x_l - x_e)/v_e$ is computed separately and exported
for plotting purposes.

\begin{figure}[H]
    \centering
    \includegraphics[width=\textwidth]{images_project3/sim_acc_controller.png}
    \caption{ACC controller block: CTG policy with minimum-set-speed
    switch and time-gap output.}
    \label{fig:sim_acc}
\end{figure}

\subsection{ISO 15622:2018 enforcement block}
Figure~\ref{fig:sim_iso} shows the post-processing chain that turns
the raw $a_\text{target}$ into a standard-compliant signal:
\begin{itemize}
    \item the leftmost switch tests $v_e > 5$\,m/s; when the ego is
          below the threshold, $a_\text{target}$ is forced to
          $\min(a_\text{target},0)$, preventing the system from
          requesting automatic positive acceleration below the
          operating-speed limit;
    \item the FIR filter implements a $2$-second moving average of
          $a_\text{target}$, used by the subsequent comparator: if the
          averaged deceleration falls below $-3.5$\,m/s$^2$, the
          switch substitutes the filtered (saturated) value for the
          raw command, preserving the standard limit;
    \item the \textbf{Jerk Rate Limiter} bounds $|\dot a|$ by
          $2.5$\,m/s$^3$;
    \item a final saturation block at $[-3.5,\,+2.0]$\,m/s$^2$ closes
          the chain.
\end{itemize}
Together these four operations cover all the ISO~15622:2018 KPIs of
Table~\ref{tab:iso}.

\begin{figure}[H]
    \centering
    \includegraphics[width=\textwidth]{images_project3/sim_ISO.png}
    \caption{ISO 15622:2018 enforcement block: minimum-speed gate,
    FIR-based deceleration limit, jerk limiter and final saturation.}
    \label{fig:sim_iso}
\end{figure}

\subsection{Plant-inversion low-level controller}
The low-level controller is implemented as the MATLAB Function
\texttt{inverse\_pedals} that returns $(h_a,h_b)$ from the triplet
$(a_\text{target},\,v,\,\tau_g)$. The traction/braking selection
follows directly from the sign of $F_\text{req}$ of~\eqref{eq:Freq}:
\begin{lstlisting}
F_res     = A + B*v^2 + C*v^4;
F_req     = me*a_target + F_res;
F_mot_max = (Te_max*eta_g*eta_t)/(Re*tau_g*tau_t);
F_brk_max = (Tb_max/Re)*(1 + 1/Kb);

if F_req > 0          % traction phase
    ha =  F_req / F_mot_max;
    hb =  0;
else                  % braking phase
    ha =  0;
    hb = -F_req / F_brk_max;
end
\end{lstlisting}
The function is purely feed-forward: it inverts the simplified
model~\eqref{eq:long_dyn} under the assumption that the lumped
parameters $(m_e,A,B,C,T_{e,\max},T_{b,\max},\ldots)$ used in the
inversion coincide with those of the full simulator. As shown in
Section~\ref{sec:res_function}, this assumption is not exactly
verified, which generates a steady-state acceleration drift.

\subsection{Augmenting the inversion with a PI controller}
\label{sec:pi}
A PI loop is wrapped around the acceleration error
$e_a = a_\text{target} - \dot v$ to compensate for the residual
plant-model mismatch. Its transfer function is
\begin{equation}
    C_\text{PI}(s) = P + \frac{I}{s},\qquad
    P = k_p = \tfrac{1}{3},\qquad
    I = \frac{k_p}{\tau_\text{virtual}} = \frac{1/3}{0.2} = \tfrac{5}{3},
    \label{eq:PI}
\end{equation}
with $\tau_\text{virtual}=0.2$\,s the desired closed-loop time
constant. The PI output is added to $a_\text{target}$ before the
inverse-pedals function: the integral channel removes the steady-state
bias, while the proportional channel keeps the transient response of
the inversion. The control parameters are summarised in
Table~\ref{tab:acc_params}.

\begin{table}[H]
    \centering
    \begin{tabular}{lcl}
        \toprule
        \textbf{Parameter} & \textbf{Value} & \textbf{Description} \\
        \midrule
        $h$                & $2.7$\,s             & Design time gap \\
        $\lambda$          & $0.8$                & CTG spacing-error gain \\
        $a_\text{max,acc}$ & $+2.0$\,m/s$^2$      & ISO max acceleration \\
        $a_\text{max,dec}$ & $-3.5$\,m/s$^2$      & ISO max deceleration \\
        $\dot a_\text{max}$& $-2.5$\,m/s$^3$      & ISO max neg. jerk \\
        $v_\text{min,acc}$ & $5$\,m/s             & ISO min op. speed \\
        $v_\text{set,min}$ & $7$\,m/s             & ISO min set speed \\
        $k_p$              & $1/3$                & PI proportional gain \\
        $\tau_\text{virtual}$ & $0.2$\,s          & PI virtual time constant \\
        \bottomrule
    \end{tabular}
    \caption{ACC and low-level controller design parameters.}
    \label{tab:acc_params}
\end{table}

%======================================================================
\section{Simulation Results}
%======================================================================
All the simulations are run on the saloon configuration (vehicle
choice 1) in dry conditions ($\mu = \mu_\text{dry}$). The ego is
initialised with $v_0 = 30$\,m/s and the leader $100$\,m ahead, in
third gear; the integration step is $T_s = 10$\,ms.

%----------------------------------------------------------------------
\subsection{Step 1 -- Plant-inversion only (\texttt{inverse\_pedals})}
\label{sec:res_function}
%----------------------------------------------------------------------
In the first simulation the PI loop is disabled and the inverse-pedals
function is used as the sole low-level controller. The lead vehicle
follows the aggressive speed profile of Figure~\ref{fig:prof2_speed}
(see Section~\ref{sec:res_prof2} for the profile itself). The ACC
results are shown in Figures~\ref{fig:func_metrics},
\ref{fig:func_leader} and~\ref{fig:func_inputs}.

\begin{figure}[H]
    \centering
    \includegraphics[width=0.95\textwidth]{images_project3/function_ACC Metrics.png}
    \caption{Step~1 (plant inversion only): ACC metrics. Top:
    commanded vs.\ actual acceleration. Middle: acceleration tracking
    error. Bottom: time gap (target $h = 2.7$\,s).}
    \label{fig:func_metrics}
\end{figure}

The top panel of Figure~\ref{fig:func_metrics} reveals the issue: the
\emph{actual} acceleration (green) stays systematically \emph{below}
the commanded one (blue) by roughly $0.3$\,m/s$^2$ throughout the run.
The middle panel quantifies the bias: the acceleration error oscillates
around $+0.2\,\div\,+0.4$\,m/s$^2$ for the entire simulation, with
positive transients up to $0.7$\,m/s$^2$ in correspondence with the
fastest speed transitions of the leader. The cause is the plant-model
mismatch between the analytical model used in the inversion
\eqref{eq:long_dyn} and the full Simulink simulator: the simplified
expression neglects, among others, the transmission-shifting
dynamics, the clutch slip, the engine-torque map non-linearities and
the wheel-slip contribution to the propulsive force. The cumulative
effect is a deterministic bias on $\dot v$ that the feed-forward
controller, by construction, cannot remove.

Despite the bias, the time gap (bottom panel) stays remarkably close
to the target $h = 2.7$\,s -- the CTG outer loop progressively
absorbs the acceleration mismatch by adjusting the spacing. The
follower position thus remains close to the desired one
(Figure~\ref{fig:func_leader}), although the speed traces show a
visible delay of the ego with respect to the leader, particularly in
the high-acceleration intervals around $t = 60$\,s and $t = 150$\,s.

\begin{figure}[H]
    \centering
    \includegraphics[width=0.95\textwidth]{images_project3/function_Leader vs Follower.png}
    \caption{Step~1 (plant inversion only): leader vs.\ ego trajectory
    (top) and speed (bottom) over the $200$\,s horizon.}
    \label{fig:func_leader}
\end{figure}

\begin{figure}[H]
    \centering
    \includegraphics[width=0.95\textwidth]{images_project3/function_Vehicle Control Inputs.png}
    \caption{Step~1 (plant inversion only): pedal commands. Brake
    $h_b$ (top) and throttle $h_a$ (bottom). The initial brake spike
    occurs while the ego decelerates from $30$\,m/s to the
    leader-matched cruise speed.}
    \label{fig:func_inputs}
\end{figure}

The pedal traces in Figure~\ref{fig:func_inputs} confirm that the
function output is well-behaved -- a short initial brake intervention
during the catch-up transient, followed by an almost-throttle-only
regime -- but the persistent acceleration bias suggests that an
integral action is required to drive the steady-state error to zero.
This motivates Step~2 of the project.

%----------------------------------------------------------------------
\subsection{Step 2 -- PI-augmented control, Profile 1 (mild)}
\label{sec:res_prof1}
%----------------------------------------------------------------------
The PI controller of equation~\eqref{eq:PI} is now enabled and the
lead vehicle plays back a mild speed profile, oscillating in the
$14\,\div\,16.5$\,m/s range (Figure~\ref{fig:prof1_speed}).

\begin{figure}[H]
    \centering
    \includegraphics[width=0.95\textwidth]{images_project3/speed_profile1.png}
    \caption{Speed profile~1: mild urban-like trace, $v_l \in
    [14,\,16.5]$\,m/s over $200$\,s.}
    \label{fig:prof1_speed}
\end{figure}

The initial transient is reported in Figure~\ref{fig:prof1_init}. The
ego starts at $30$\,m/s while the leader cruises near $15$\,m/s, so
the ACC must decelerate aggressively. The commanded acceleration
saturates at $-3.5$\,m/s$^2$ for about $5$\,s (top panel), exactly the
ISO~15622 lower bound, and is held there by the FIR-filter switch of
the ISO block. After the speed catches up, around $t = 8$\,s, the
commanded acceleration jumps to $+2$\,m/s$^2$ -- again the ISO upper
limit -- to recover the time gap, which had collapsed to $\sim 1.3$\,s
(bottom panel) during the catch-up. The PI loop is unable to fully
follow the saturated commands during the transient: the acceleration
error peaks at $5$\,m/s$^2$ around $t = 9$\,s when the commanded
acceleration changes sign. The transient is over by $t \approx 25$\,s
and the time gap converges to the target $2.7$\,s.

\begin{figure}[H]
    \centering
    \includegraphics[width=0.95\textwidth]{images_project3/prof1_beginning_metrics.png}
    \caption{Profile~1, transient phase ($0\le t\le 31$\,s):
    saturation of $a_\text{target}$ at the ISO limits, large
    acceleration error during the sign reversal at $t\approx 9$\,s
    and time-gap recovery to $2.7$\,s.}
    \label{fig:prof1_init}
\end{figure}

In the steady-state regime, the PI action delivers practically perfect
acceleration tracking. Figure~\ref{fig:prof1_ss} zooms on a one-second
window around $t = 70$\,s: the acceleration error is bounded within
$\pm 0.02$\,m/s$^2$ -- one order of magnitude smaller than the residual
drift of Step~1 -- and the time gap is indistinguishable from the
target trace. This is the key qualitative result of the project: the
integral channel of the PI eliminates the steady-state bias that the
pure inversion of Step~1 could not handle.

\begin{figure}[H]
    \centering
    \includegraphics[width=0.95\textwidth]{images_project3/prof1_steady-state ACC Metrics.png}
    \caption{Profile~1, steady-state window ($t\approx 70$\,s):
    acceleration error confined within $\pm 0.02$\,m/s$^2$ and time
    gap perfectly aligned with the target $h = 2.7$\,s.}
    \label{fig:prof1_ss}
\end{figure}

The leader-vs.-follower comparison in Figure~\ref{fig:prof1_leader} is
restricted to the steady-state portion of the run ($t\ge 50$\,s) to
avoid the spurious large initial gap. The ego speed (red) tracks the
leader speed (cyan) with a barely visible delay; the position curves
are essentially parallel, with the constant spacing dictated by the
CTG policy.

\begin{figure}[H]
    \centering
    \includegraphics[width=0.95\textwidth]{images_project3/prof1_Leader vs Follower.png}
    \caption{Profile~1, $t\ge 50$\,s: ego-vs.-leader position (top)
    and speed (bottom). The constant-headway separation is maintained
    throughout.}
    \label{fig:prof1_leader}
\end{figure}

\begin{figure}[H]
    \centering
    \includegraphics[width=0.95\textwidth]{images_project3/prof1_Vehicle Control Inputs.png}
    \caption{Profile~1: pedal commands. A single full-brake event at
    $t\approx 3$\,s ($h_b\to 1$) is followed by an almost-throttle-only
    operation, with $h_a$ oscillating around $0.3$.}
    \label{fig:prof1_inputs}
\end{figure}

The pedal traces (Figure~\ref{fig:prof1_inputs}) confirm a comfortable
operating regime: after the initial full-brake event the throttle
oscillates between $0.2$ and $0.5$, well inside the actuator range,
and the brake stays at zero except for occasional micro-pulses at
$t\approx 22$\,s.

%----------------------------------------------------------------------
\subsection{Step 2 -- PI-augmented control, Profile 2 (aggressive)}
\label{sec:res_prof2}
%----------------------------------------------------------------------
The second test exercises a more demanding scenario: the leader speed
ramps from $\sim 15$\,m/s to $\sim 28$\,m/s and back, covering both
urban and extra-urban regimes (Figure~\ref{fig:prof2_speed}).

\begin{figure}[H]
    \centering
    \includegraphics[width=0.95\textwidth]{images_project3/speed_profile2.png}
    \caption{Speed profile~2: aggressive trace, $v_l \in
    [13,\,28]$\,m/s over $200$\,s.}
    \label{fig:prof2_speed}
\end{figure}

Figure~\ref{fig:prof2_metrics} reports the ACC metrics on the
steady-state window. The commanded acceleration (blue) oscillates in
the $-0.5\,\div\,+1.0$\,m/s$^2$ range, well inside the ISO bounds,
following the speed fluctuations of the leader. The actual
acceleration (green) sticks to the command with a worst-case error
$|e_a|\le 0.3$\,m/s$^2$, and the steady error is centred on zero --
i.e.\ no bias, in contrast with Step~1. The time gap stays within
$\pm 0.05$\,s of the target, with the largest deviations occurring at
the strongest acceleration ramps (around $t = 70$\,s and $t = 88$\,s).

\begin{figure}[H]
    \centering
    \includegraphics[width=0.95\textwidth]{images_project3/prof2_ACC Metrics.png}
    \caption{Profile~2, steady-state window ($65\le t\le 140$\,s):
    acceleration tracking and time-gap regulation under the
    aggressive leader profile.}
    \label{fig:prof2_metrics}
\end{figure}

The leader-vs.-follower comparison (Figure~\ref{fig:prof2_leader})
shows that the ego closely follows the leader through the whole speed
envelope, with a marginally larger phase delay during the high-speed
plateau ($120\,\div\,160$\,s, $v_l\approx 27$\,m/s). The pedal traces
(Figure~\ref{fig:prof2_inputs}) reflect the higher dynamic demand:
several throttle saturations at $h_a = 1$ occur during the
acceleration phases (the controller asks for $+2$\,m/s$^2$ but the
engine cannot deliver more in the current gear), while the brake
intervenes only sporadically with peaks below $0.05$.

\begin{figure}[H]
    \centering
    \includegraphics[width=0.95\textwidth]{images_project3/prof2_Leader vs Follower.png}
    \caption{Profile~2: ego-vs.-leader position (top) and speed
    (bottom). The vehicle tracks the leader through the full
    $13\,\div\,28$\,m/s range.}
    \label{fig:prof2_leader}
\end{figure}

\begin{figure}[H]
    \centering
    \includegraphics[width=0.95\textwidth]{images_project3/prof2_Vehicle Control Inputs.png}
    \caption{Profile~2: pedal commands. The throttle saturates at
    $h_a = 1$ during the acceleration phases, while the brake stays
    near zero apart from short micro-pulses.}
    \label{fig:prof2_inputs}
\end{figure}

The aggressive scenario thus confirms the conclusions drawn on
Profile~1: the PI-augmented controller delivers zero-bias acceleration
tracking and tight headway regulation across a wide speed range, all
while respecting the ISO~15622:2018 KPIs enforced by the dedicated
Simulink block.

%======================================================================
\section{Benchmark}
\label{sec:bench_p3}
%======================================================================
Table~\ref{tab:bench_p3} summarises the quantitative comparison of the
three simulated scenarios. The figures are read directly from the
time histories of $e_a(t)$, $t_\text{gap}(t)$ and the pedal commands,
with the steady-state window starting at $t\ge 50$\,s for all the
profiles. ``ISO'' columns report whether each KPI of
Table~\ref{tab:iso} is satisfied throughout the simulation.

\begin{table}[H]
    \centering
    \footnotesize
    \begin{tabular}{lccccc}
        \toprule
        \textbf{Scenario} &
        \textbf{$|e_a|_\text{ss}$} &
        \textbf{$|t_\text{gap}-h|_\text{ss}$} &
        \textbf{Max $|\dot v|$} &
        \textbf{ISO compl.} &
        \textbf{Behaviour} \\
        \midrule
        Step~1 -- inversion only, aggr.\ prof.\ &
            $0.20\,\div\,0.40$\,m/s$^2$ & $<0.1$\,s & $\le 2.0$\,m/s$^2$ &
            yes & biased, $\dot v < a_\text{target}$ \\
        Step~2 -- PI, Profile~1 (mild) &
            $<0.02$\,m/s$^2$ & $<0.02$\,s & $\le 3.5$\,m/s$^2$ &
            yes & best tracking, no bias \\
        Step~2 -- PI, Profile~2 (aggressive) &
            $<0.30$\,m/s$^2$ & $<0.05$\,s & $\le 2.0$\,m/s$^2$ &
            yes & robust, throttle saturated in $\dot v_+$ ramps \\
        \bottomrule
    \end{tabular}
    \caption{Benchmark of the three ACC simulations. Steady-state
    figures are evaluated for $t\ge 50$\,s; the transient catch-up
    region is excluded.}
    \label{tab:bench_p3}
\end{table}

Three observations stand out:
\begin{itemize}
    \item The pure plant-inversion controller of Step~1 is
          intrinsically biased. The acceleration error of $\approx
          0.3$\,m/s$^2$ stems from the mismatch between the simplified
          model used in the inversion and the full simulator
          (transmission, clutch, engine-map non-linearities) and
          cannot be eliminated by retuning the feed-forward gains.
    \item Adding a PI loop on the acceleration error -- with the
          virtual-time-constant tuning $P = k_p,\,I = k_p/\tau_v$ --
          completely removes the steady-state drift and reduces the
          tracking error by more than one order of magnitude.
    \item Compliance with ISO~15622:2018 is guaranteed by the
          dedicated post-processing block (minimum-speed gate, FIR
          deceleration filter, jerk limiter, final saturation) in
          every scenario, including the most demanding accelerations
          of Profile~2 where the upstream CTG law would otherwise
          ask for $|\dot v|>2$\,m/s$^2$.
\end{itemize}

%======================================================================
\section{Conclusions}
%======================================================================
The project addressed the design of a full ACC stack on a nonlinear
longitudinal vehicle simulator, starting from the high-level
constant-time-gap law down to the pedal-level commands.
\begin{itemize}
    \item \textbf{High-level controller.} A classical CTG PD law on
          the spacing error, with $h = 2.7$\,s and $\lambda = 0.8$,
          regulates the inter-vehicle distance and delivers a target
          acceleration that respects, by design, the steady-state
          headway requirement.
    \item \textbf{Low-level controller -- Step 1.} The analytical
          inversion of the simplified longitudinal model is
          straightforward and computationally cheap, but
          plant-simulator mismatch generates a persistent $\approx
          0.3$\,m/s$^2$ acceleration bias. The follower lags slightly
          behind the leader and the closed-loop CTG action absorbs the
          bias only through a small permanent spacing offset.
    \item \textbf{Low-level controller -- Step 2.} Augmenting the
          inversion with a PI loop on the acceleration error (tuned
          via the virtual-time-constant rule, $\tau_v = 0.2$\,s)
          removes the bias entirely. On the mild Profile~1 the
          steady-state acceleration error drops below
          $0.02$\,m/s$^2$, and on the aggressive Profile~2 it stays
          below $0.30$\,m/s$^2$ even during the strongest
          acceleration ramps, with the time gap permanently within
          $\pm 0.05$\,s of the target.
    \item \textbf{ISO 15622:2018 compliance.} A dedicated block
          combining a minimum-speed gate, a $2$-second FIR filter on
          the deceleration, a jerk rate limiter at
          $2.5$\,m/s$^3$ and a final saturation at
          $[-3.5,\,+2.0]$\,m/s$^2$ enforces all the five KPIs of the
          standard, both in the transient and the steady-state
          regions of the simulations.
\end{itemize}

The key methodological takeaway is that \emph{plant inversion alone
is insufficient when a non-negligible plant-model mismatch is present}:
the feed-forward action handles the bulk of the conversion from
$a_\text{target}$ to $(h_a,h_b)$, but a complementary feedback term --
here an integral action driven by the acceleration error -- is
indispensable to achieve zero steady-state drift. A natural extension
of this work is the introduction of an anti-windup logic on the PI
integrator, which would speed up the recovery from the saturated
acceleration commands observed at the beginning of every simulation
(see Figure~\ref{fig:prof1_init}), and the adoption of a gain-scheduled
PI whose proportional gain depends on the engaged gear, to better
match the gear-dependent throttle authority that the current
single-gain design only approximates.
